% ===================================================================
% Section 4: PI + prepínač pracovných bodov 
% ===================================================================

\section{PI regulácia s prepínaním pracovných bodov}
\label{sec:pi_switching}

% -------------------------------------------------------------------
% Subsection 4.1 - Prepínač pracovných bodov
% -------------------------------------------------------------------

\subsection{Prepínač pracovných bodov}

V tejto časti sa budeme zaoberať návrhom a testovaním prepínača pracovných bodov.  
Prepínač bude predstavovať pomalý regulátor, ktorého úlohou bude zabezpečiť robustné prechody medzi rôznymi pracovnými bodmi systému TS05.  
Cieľom je vytvoriť taký spôsob riadenia, ktorý umožní spoľahlivý presun z ľubovoľného pracovného bodu do ľubovoľného iného pracovného bodu bez neželaných dynamických javov, ako je presýtenie akčného zásahu alebo príliš veľké prekmitové správanie.

Okrem samotného prepínača bola v tejto časti rozšírená aj pracovná oblasť regulácie o dva ďalšie pracovné body systému – pre napätia špirály 2~V a 6~V. Pre každý z týchto pracovných bodov bol postup opakovaný presne podľa postupu uvedeného v sekcii~\ref{sec:identifikacia},  
teda boli:

\begin{itemize}
    \item vykonané merania reakcií systému na jednotkové skoky $\pm 0.5$~V okolo pracovného bodu,
    \item z nameraných údajov bola identifikovaná prenosová funkcia v danom pracovnom bode,
    \item pre každý pracovný bod bol navrhnutý PI regulátor s vlastnými parametrami,
    \item boli určené hodnoty $y_{pb}$ a rozsah výstupu pre $\pm 0.5$~V zmenu vstupu.
\end{itemize}

Prehľad všetkých troch pracovných bodov je uvedený v tabuľke~\ref{tab:pb_overview},  
kde sú uvedené identifikované prenosové funkcie, hodnoty pracovného bodu $y_{pb}$ a parametre jednotlivých PI regulátorov.

\begin{table}[h!]
    \centering
    \caption{Parametre pracovných bodov, identifikované prenosové funkcie a PI regulátory}
    \label{tab:pb_overview}
    \begin{tabular}{|c|c|c|c|c|}
        \hline
        \textbf{$u_{pb}$} & \textbf{Prenosová funkcia} & \textbf{$y_{pb}$} & \textbf{$\Delta y_{pb}$ pri $u_{pb}$ $\pm 0.5$~V} & \textbf{PI parametre} \\ \hline

        2~V &
        $\displaystyle \frac{5.79}{s^2 + 44.21 s + 7.276}$ &
        4.4884 &
        $\pm 0.41$ &
        $K_p = 2.77$, $K_i = 1.48$ \\ \hline

        4~V &
        $\displaystyle \frac{6.074}{s^2 + 59.01 s + 8.95}$ &
        6.3412 &
        $\pm 0.35$ &
        $K_p = 3.58$, $K_i = 1.84$ \\ \hline

        6~V &
        $\displaystyle \frac{7.161}{s^2 + 52.87 s + 21.37}$ &
        7.7441 &
        $\pm 0.18$ &
        $K_p = 1.83$, $K_i = 2.20$ \\ \hline
    \end{tabular}
\end{table}

Keďže pôvodné statické charakteristiky systému boli merané ešte počas letného obdobia, bolo rozhodnuté hodnoty $y_{pb}$ pre tieto experimenty zmerať znova, rovnakým spôsobom ako v sekcii~\ref{sec:staticke_vlastnosti}.  
V tabuľke sú preto uvedené aktuálne hodnoty pracovných bodov, ktoré vernejšie reprezentujú stav systému počas meraní v zimnom období.


% -------------------------------------------------------------------
% Subsubsection 4.1.1 - Opis a schéma riadenia s použitím prepínača
% -------------------------------------------------------------------

\subsubsection{Opis a schéma riadenia s použitím prepínača}

V tejto časti je uvažovaná riadiaca štruktúra,
ktorá pozostáva z dvoch PI regulátorov.
Prvý z nich predstavuje globálny PI regulátor,
ktorý pracuje s absolútnymi veličinami systému.
Druhý predstavuje lokálny PI regulátor,
ktorý pracuje s normalizovanými veličinami
v okolí aktuálneho pracovného bodu.

Lokálny regulátor využíva normalizované premenné

\begin{equation}
y_N = y - y_{pb},
\qquad
u_N = u - u_{pb},
\label{eq:pi_switch_normalized_variables}
\end{equation}

kde $y_{pb}$ a $u_{pb}$
predstavujú výstupnú a vstupnú hodnotu príslušného pracovného bodu.
Parametre lokálneho PI regulátora
sú volené pomocou princípu gain scheduling,
pričom sa používajú hodnoty uvedené v tabuľke~\ref{tab:pb_overview}.

Globálny regulátor pracuje s absolútnou žiadanou hodnotou.
Aby prechod medzi pracovnými bodmi neprebiehal príliš prudko,
globálna žiadaná hodnota je pred vstupom do regulátora filtrovaná
prenosom druhého rádu v tvare

\begin{equation}
G_f(s)=\frac{1}{(10s+1)^2}.
\label{eq:pi_switch_global_filter_tf}
\end{equation}

Takto zvolený filter zabezpečuje dostatočne plynulý priebeh
globálnej žiadanej hodnoty,
v dôsledku čoho prechod medzi pracovnými bodmi
prebieha pomalšie a s menším rizikom vzniku neželaných prechodových javov.

O tom, ktorý PI regulátor bude v danom okamihu aktívny,
rozhoduje blok logiky prepínača.
Jeho úloha je dvojitá.
Po prvé určuje,
či má byť aktívny globálny alebo lokálny regulátor.
Po druhé spracúva globálnu žiadanú hodnotu
a vytvára jej filtrovanú verziu pre globálnu reguláciu.

Navrhnutá bloková schéma algoritmu
je znázornená na obr.~\ref{fig:pi_switch_logic_scheme}.

\begin{center}
    \vbox{
        \makebox[\textwidth][c]{%
            \includegraphics[width=1.2\textwidth]{SWITCH_block_logic_scheme.jpg}
        }
        \figcaption{Bloková schéma riadenia s globálnym a lokálnym PI regulátorom}
        \label{fig:pi_switch_logic_scheme}
    }
\end{center}

Logika prepínania je založená na nasledujúcom princípe.
Ak je detegovaná skoková zmena globálnej žiadanej hodnoty,
aktivuje sa globálny PI regulátor,
ktorý zabezpečuje prechod medzi pracovnými bodmi.
Ak globálna regulácia nie je aktívna,
systém sa automaticky prepne do režimu lokálnej regulácie,
v ktorom sa používa PI regulátor
naladený pre aktuálny pracovný bod.

V ďalšej časti budú predstavené dve konkrétne realizácie tejto logiky.
Budú sa líšiť predovšetkým spôsobom rozhodovania o tom,
kedy sa globálna regulácia ukončí
a kedy sa aktivuje lokálna regulácia.

\paragraph{Sledovanie integrátora}

Dôležitou požiadavkou navrhnutej štruktúry
je zabezpečenie plynulého prepínania medzi globálnym
a lokálnym PI regulátorom.
Ak by sa po deaktivácii jedného regulátora
jeho integračná zložka nechala bez aktualizácie,
po opätovnom aktivovaní by mohla mať hodnotu,
ktorá nezodpovedá aktuálne realizovanému akčnému zásahu systému.
To by viedlo ku skokovej zmene riadenia
a k neželanej dynamike pri prepnutí.

Z tohto dôvodu je pri neaktívnom regulátore
použitý jednoduchý sledovací mechanizmus integrátora.
Jeho úlohou je priebežne nastavovať integračný stav tak,
aby po aktivovaní regulátora nedošlo
k skokovej zmene jeho výstupu.
V prípade lokálneho regulátora sa pri tomto výpočte
využíva aj hodnota $u_{pb}$,
ktorá je počas globálnej regulácie zvolená
podľa aktuálne uvažovaného pracovného bodu
z tabuľky~\ref{tab:pb_overview}.
Takto sa dosiahne plynulé prepnutie medzi oboma regulačnými vetvami.

\paragraph{Zákon riadenia}

Označme binárne aktivačné signály lokálneho a globálneho regulátora ako
$a_l$ a $a_g$, pričom platí

\begin{equation}
a_l, a_g \in \{0,1\},
\qquad
a_l + a_g = 1.
\label{eq:pi_switch_activity_condition}
\end{equation}

Chyba lokálneho regulátora je definovaná vzťahom

\begin{equation}
e_l = r_l - \left( y - y_{pb} \right),
\label{eq:pi_switch_local_error}
\end{equation}

kde $r_l$ predstavuje lokálnu žiadanú hodnotu.

Filtrovaná globálna žiadaná hodnota je označená ako $r_{g,f}$,
pričom chyba globálneho regulátora je

\begin{equation}
e_g = r_{g,f} - y.
\label{eq:pi_switch_global_error}
\end{equation}

Lokálny PI regulátor má tvar

\begin{equation}
u_l = \left( u_{pb} + K_{p,l} e_l + K_{i,l} I_l \right) a_l,
\label{eq:pi_switch_local_control}
\end{equation}

a globálny PI regulátor je definovaný vzťahom

\begin{equation}
u_g = \left( K_{p,g} e_g + K_{i,g} I_g \right) a_g.
\label{eq:pi_switch_global_control}
\end{equation}

Celkový akčný zásah je daný súčtom oboch vetiev

\begin{equation}
u = u_l + u_g.
\label{eq:pi_switch_total_control}
\end{equation}

Po dosadení \eqref{eq:pi_switch_local_control}
a \eqref{eq:pi_switch_global_control}
do \eqref{eq:pi_switch_total_control}
dostávame rozvinutý tvar zákona riadenia

\begin{equation}
u =
\left( u_{pb} + K_{p,l} e_l + K_{i,l} I_l \right) a_l
+
\left( K_{p,g} e_g + K_{i,g} I_g \right) a_g.
\label{eq:pi_switch_total_control_expanded}
\end{equation}

Keďže v každom okamihu je aktívna práve jedna regulačná vetva,
možno zákon riadenia zapísať aj v tvare

\begin{equation}
u =
\begin{cases}
u_{pb} + K_{p,l} e_l + K_{i,l} I_l, & a_l = 1,\; a_g = 0, \\[1mm]
K_{p,g} e_g + K_{i,g} I_g, & a_l = 0,\; a_g = 1.
\end{cases}
\label{eq:pi_switch_total_control_cases}
\end{equation}

Pri aktívnom regulátore sa integračný stav aktualizuje štandardne
na základe regulačnej chyby.
Pri neaktívnom regulátore sa integrátor priamo dopočíta
z aktuálneho akčného zásahu,
aby bol pripravený na plynulé prevzatie riadenia.

Pre lokálny regulátor sa pri neaktívnom stave používa vzťah

\begin{equation}
I_l = \frac{(u-u_{pb}) - K_{p,l} e_l}{K_{i,l}},
\qquad a_l = 0,
\label{eq:pi_switch_local_tracking}
\end{equation}

a pre globálny regulátor analogicky

\begin{equation}
I_g = \frac{u - K_{p,g} e_g}{K_{i,g}},
\qquad a_g = 0.
\label{eq:pi_switch_global_tracking}
\end{equation}

Takto definovaný spôsob sledovania integrátora
zabezpečuje,
že pri prepnutí medzi lokálnou a globálnou regulačnou vetvou
nedochádza k skokovej zmene akčného zásahu.

% -------------------------------------------------------------------
% Subsubsection 4.1.2 - Globálna regulácia po určitý čas
% -------------------------------------------------------------------

\subsubsection{Globálna regulácia po určitý čas}
\label{sec:pi_switch_timer_based}

Prvou uvažovanou realizáciou prepínacej logiky
je regulácia založená na pevne zvolenom čase pôsobenia globálneho PI regulátora.
Predpokladá sa pritom,
že na prechod z jedného pracovného bodu do druhého
postačuje približne konštantný čas.
Na základe experimentálnych pozorovaní
bol tento čas určený ako

\begin{equation}
T_{sw} = 80~\mathrm{s}.
\label{eq:pi_switch_timer}
\end{equation}

Pri skokovej zmene globálnej žiadanej hodnoty
sa teda aktivuje globálny PI regulátor,
ktorý zostáva aktívny počas intervalu $T_{sw}$.
Po uplynutí tohto času sa globálna vetva deaktivuje
a riadenie automaticky preberá lokálny PI regulátor.
Pri lokálnej regulácii sa používajú hodnoty
$u_{pb}$, $y_{pb}$ a parametre PI regulátora
z tabuľky~\ref{tab:pb_overview}.

Filter globálnej žiadanej hodnoty
z rovnice~\eqref{eq:pi_switch_global_filter_tf}
je pre potreby diskrétnej implementácie
vhodné previesť do stavového priestoru.
Pre prenos

\begin{equation}
G_f(s)=\frac{R_{g,f}(s)}{R_g(s)}
=
\frac{1}{(10s+1)^2}
=
\frac{0.01}{s^2 + 0.2s + 0.01}
\label{eq:pi_switch_filter_tf_repeat}
\end{equation}

zvolíme stavový vektor

\begin{equation}
x_f(t)=
\begin{bmatrix}
x_{f1}(t)\\
x_{f2}(t)
\end{bmatrix}.
\label{eq:pi_switch_filter_state_vector}
\end{equation}

V kanonickom tvare má filter podobu

\begin{align}
\dot{x}_f(t) &= A_f x_f(t) + b_f r_g(t),
\label{eq:pi_switch_filter_state_model}
\\
r_{g,f}(t) &= c_f^T x_f(t),
\label{eq:pi_switch_filter_output_model}
\end{align}

kde

\begin{equation}
A_f =
\begin{bmatrix}
0 & 1\\
-0.01 & -0.2
\end{bmatrix},
\qquad
b_f =
\begin{bmatrix}
0\\
1
\end{bmatrix},
\qquad
c_f =
\begin{bmatrix}
0.01\\
0
\end{bmatrix}.
\label{eq:pi_switch_filter_matrices}
\end{equation}

Pri Eulerovej diskretizácii so vzorkovacou periódou $T_s$
má stavová rovnica tvar

\begin{equation}
x_f(k+1)=x_f(k)+T_s\left(A_f x_f(k)+b_f r_g(k)\right),
\label{eq:pi_switch_filter_discrete_state}
\end{equation}

a výstup filtra je určený vzťahom

\begin{equation}
r_{g,f}(k)=c_f^T x_f(k).
\label{eq:pi_switch_filter_discrete_output}
\end{equation}

Z implementačného hľadiska je preto potrebné
prenášať medzi jednotlivými krokmi
stav filtra globálnej žiadanej hodnoty,
integračné stavy oboch PI regulátorov,
predchádzajúcu hodnotu globálnej žiadanej hodnoty
a v prípade časovo riadeného prepínania aj hodnotu časovača.
Okrem toho je potrebné mať k dispozícii tabuľkové hodnoty
pracovných bodov $y_{pb}$, $u_{pb}$
a príslušné parametre lokálnych PI regulátorov,
keďže lokálna regulačná vetva je realizovaná
s využitím gain scheduling.

Pseudokód tejto realizácie je uvedený nižšie.

\begin{lstlisting}[language=Matlab,
caption={Pseudokód prepínania s globálnou reguláciou po určitý čas},
label={lst:pi_switch_timer_based}]
% Inputs:
% r_g(k) - global reference
% r_l(k) - local reference
% y(k)   - plant output
%
% States:
% x_f(2x1)    - state of global reference filter
% I_l         - integral state of local PI controller
% I_g         - integral state of global PI controller
% timer_g     - elapsed time of global mode
% r_g_prev    - previous global reference
% active_g    - activity flag of global controller
%
% Parameters:
% Ts, T_sw
% Kp_g, Ki_g  - global PI parameters
%
% Tables:
% y_pb_table, u_pb_table
% Kp_l_table, Ki_l_table

% Filter matrices
A_f = [ 0,    1;
       -0.01, -0.2 ];

b_f = [ 0;
        1 ];

c_f = [ 0.01, 0 ];

% ------------------------------------------------------------------
% 1. Detect step change of global reference
% ------------------------------------------------------------------
if r_g ~= r_g_prev
    active_g = 1;
    timer_g = 0;
end

% ------------------------------------------------------------------
% 2. Filter global reference
% ------------------------------------------------------------------
r_gf = c_f * x_f;
x_f_next = x_f + Ts * ( A_f * x_f + b_f * r_g );

% ------------------------------------------------------------------
% 3. Select local operating point and PI parameters from table
% ------------------------------------------------------------------
select y_pb, u_pb, Kp_l, Ki_l from table according to current operating point

% ------------------------------------------------------------------
% 4. Switching logic
% ------------------------------------------------------------------
if active_g == 1
    timer_g = timer_g + Ts;
    if timer_g >= T_sw
        active_g = 0;
    end
end

active_l = 1 - active_g;

% ------------------------------------------------------------------
% 5. Control errors
% ------------------------------------------------------------------
e_l = r_l - ( y - y_pb );
e_g = r_gf - y;

% ------------------------------------------------------------------
% 6. Control law
% ------------------------------------------------------------------
u_l = active_l * ( u_pb + Kp_l * e_l + Ki_l * I_l );
u_g = active_g * ( Kp_g * e_g + Ki_g * I_g );
u = u_l + u_g;

% ------------------------------------------------------------------
% 7. Integral states update
% ------------------------------------------------------------------
if active_l == 1
    I_l_next = I_l + Ts * e_l;
else
    I_l_next = ( (u - u_pb) - Kp_l * e_l ) / Ki_l;
end

if active_g == 1
    I_g_next = I_g + Ts * e_g;
else
    I_g_next = ( u - Kp_g * e_g ) / Ki_g;
end

% ------------------------------------------------------------------
% 8. Memory update
% ------------------------------------------------------------------
r_g_prev_next = r_g;
timer_g_next = timer_g;
\end{lstlisting}

Uvedený spôsob prepínania je založený na predpoklade,
že globálny PI regulátor potrebuje na presun systému
do nového pracovného bodu približne konštantný čas.
Takýto prístup je implementačne jednoduchý,
avšak jeho nevýhodou je,
že dynamika systému nie je v reálnych podmienkach úplne konštantná.
Mení sa okrem iného aj v závislosti od aktuálnych podmienok okolia,
napríklad od teploty v miestnosti.
To znamená,
že hodnoty $u_{pb}$ a $y_{pb}$ uvedené v tabuľke~\ref{tab:pb_overview}
nemusia v danom experimente presne zodpovedať skutočnému ustálenému stavu systému.

Rovnako nie je zaručené,
že za pevne stanovený čas $80$~s
sa systém vždy ustáli dostatočne blízko požadovaného pracovného bodu.
Z tohto dôvodu bol navrhnutý aj druhý spôsob prepínania,
ktorý je založený priamo na detekcii ustálenia.

% -------------------------------------------------------------------
% Subsubsection 4.1.3 - Prepínanie na základe detekcie ustálenia
% -------------------------------------------------------------------

\subsubsection{Prepínanie na základe detekcie ustálenia}
\label{sec:pi_switch_settling_detection}

Druhá uvažovaná realizácia prepínacej logiky
nie je založená na predpoklade,
že ustálenie nastane po vopred známom čase.
Namiesto toho sa globálna regulácia ukončí až vtedy,
keď sa na základe meraných údajov zistí,
že systém sa skutočne ustálil.

V tejto štruktúre má globálny PI regulátor úlohu
pomalého a robustného regulátora,
ktorý zabezpečuje prechod medzi pracovnými bodmi
v rámci celého pracovného rozsahu systému.
Lokálny PI regulátor má naopak za úlohu
presnú reguláciu v okolí aktuálneho pracovného bodu,
teda v oblasti približne
$u_{pb} \pm 0.5$~V.
Z tohto dôvodu pracuje s normalizovanými veličinami
a jeho parametre sa volia pomocou gain scheduling
na základe tabuľkových hodnôt pracovných bodov.

Pri tejto metóde sa okrem stavu filtra,
integračných stavov a predchádzajúcej hodnoty globálnej žiadanej hodnoty
uchováva aj posledných $30$ hodnôt výstupu $y$
a posledných $30$ hodnôt vstupu $u$.
Globálna regulácia sa opäť aktivuje pri skoku globálnej žiadanej hodnoty,
avšak jej deaktivácia nastane až po splnení podmienky ustálenia.

Za ustálený stav sa v tomto prípade považuje situácia,
keď sa aspoň $25$ z posledných $30$ hodnôt výstupu
nachádza v intervale

\begin{equation}
\left\langle r_g - 0.07\,r_g,\; r_g + 0.07\,r_g \right\rangle.
\label{eq:pi_switch_settling_band}
\end{equation}

Po splnení tejto podmienky
sa ako aktuálna hodnota pracovného bodu určí aritmetický priemer
posledných $30$ hodnôt výstupu a vstupu, teda

\begin{equation}
y_{pb}^{\ast} = \frac{1}{30}\sum_{i=1}^{30} y_i,
\qquad
u_{pb}^{\ast} = \frac{1}{30}\sum_{i=1}^{30} u_i.
\label{eq:pi_switch_estimated_operating_point}
\end{equation}

Takto získané hodnoty reprezentujú odhad skutočného pracovného bodu,
ktorý bol dosiahnutý globálnou reguláciou.
Následne sa vyberú parametre lokálneho PI regulátora
podľa toho,
ku ktorému pracovnému bodu z tabuľky~\ref{tab:pb_overview}
je hodnota $y_{pb}^{\ast}$ najbližšia.

Pseudokód tejto realizácie je uvedený nižšie.

\begin{lstlisting}[language=Matlab,
caption={Pseudokód prepínania na základe detekcie ustálenia},
label={lst:pi_switch_settling_based}]
% Inputs:
% r_g(k) - global reference
% r_l(k) - local reference
% y(k)   - plant output
%
% States:
% x_f(2x1)    - state of global reference filter
% I_l         - integral state of local PI controller
% I_g         - integral state of global PI controller
% y_hist(30)  - history of output values
% u_hist(30)  - history of input values
% r_g_prev    - previous global reference
% active_g    - activity flag of global controller
%
% Parameters:
% Ts
% Kp_g, Ki_g  - global PI parameters
%
% Tables:
% y_pb_table, u_pb_table
% Kp_l_table, Ki_l_table

% Filter matrices
A_f = [ 0,    1;
       -0.01, -0.2 ];

b_f = [ 0;
        1 ];

c_f = [ 0.01, 0 ];

% ------------------------------------------------------------------
% 1. Detect step change of global reference
% ------------------------------------------------------------------
if r_g ~= r_g_prev
    active_g = 1;
end

% ------------------------------------------------------------------
% 2. Filter global reference
% ------------------------------------------------------------------
r_gf = c_f * x_f;
x_f_next = x_f + Ts * ( A_f * x_f + b_f * r_g );

% ------------------------------------------------------------------
% 3. Update histories
% ------------------------------------------------------------------
shift y_hist and append y(k)
shift u_hist and append u(k)

% ------------------------------------------------------------------
% 4. Select local operating point and PI parameters from table
% ------------------------------------------------------------------
select y_pb, u_pb, Kp_l, Ki_l from table according to current operating point

% ------------------------------------------------------------------
% 5. Control errors
% ------------------------------------------------------------------
e_l = r_l - ( y - y_pb );
e_g = r_gf - y;

% ------------------------------------------------------------------
% 6. Control law
% ------------------------------------------------------------------
u_l = (1 - active_g) * ( u_pb + Kp_l * e_l + Ki_l * I_l );
u_g = active_g * ( Kp_g * e_g + Ki_g * I_g );
u = u_l + u_g;

% ------------------------------------------------------------------
% 7. Settling detection
% ------------------------------------------------------------------
count_ok = number of elements in y_hist within r_g +/- 7%

if active_g == 1 && count_ok >= 25
    y_pb_est = mean( y_hist );
    u_pb_est = mean( u_hist );

    select local PI parameters corresponding to nearest tabulated y_pb
    set y_pb = y_pb_est;
    set u_pb = u_pb_est;

    active_g = 0;
end

active_l = 1 - active_g;

% ------------------------------------------------------------------
% 8. Integral states update
% ------------------------------------------------------------------
if active_l == 1
    I_l_next = I_l + Ts * e_l;
else
    I_l_next = ( (u - u_pb) - Kp_l * e_l ) / Ki_l;
end

if active_g == 1
    I_g_next = I_g + Ts * e_g;
else
    I_g_next = ( u - Kp_g * e_g ) / Ki_g;
end

% ------------------------------------------------------------------
% 9. Memory update
% ------------------------------------------------------------------
r_g_prev_next = r_g;
\end{lstlisting}

Táto metóda je z výpočtového hľadiska zložitejšia,
pretože vyžaduje uchovávanie histórie signálov
a priebežné vyhodnocovanie podmienky ustálenia.
Na druhej strane je však podstatne robustnejšia,
pretože nepredpokladá fixný čas ustálenia,
ale vychádza zo skutočného správania systému.

Keďže perióda regulácie je $T_s = 0.1$~s,
čo predstavuje z pohľadu použitého hardvéru
dostatočne veľký čas na vykonanie potrebných výpočtov,
bolo rozhodnuté implementovať práve tento prístup.
V porovnaní s časovo riadeným prepínaním
poskytuje spoľahlivejšie ukončenie globálnej regulácie
a lepšie zohľadňuje zmeny dynamiky systému
spôsobené reálnymi prevádzkovými podmienkami.